% ===========================================================================
\section{Sets of curvilinear convergence}\label{sec:cc}
% ===========================================================================
Sources: \cite{MMJ}; \cite[\S4, Lemmas 7--8, Theorem 3]{Diss}. We work in $\D$ and transfer to $\Hp$.

\begin{definition}[accessible points and crosscut arcs]\label{def:acc}
\lean{TJK.McMillan.Gamma, TJK.McMillan.crossArcOpen, TJK.McMillan.I0, TJK.McMillan.acc}
\uses{def:arcs, fact:crosscut}
For a crosscut $\gamma$ of $\D$ with $0\notin\gamma$, let $U$ be the component of $\D\setminus\gamma$ containing
$0$ and $L(\gamma)^*=C\setminus\overline U$; by Theorem~\ref{fact:crosscut} it is the open arc of $C$ between the
endpoints of $\gamma$ on the side of the other component. For a domain $\Omega\subseteq\D\setminus\{0\}$ let
$\Gamma(\Omega)$ be the set of crosscuts $\gamma$ with $\gamma\cap\D\subseteq\Omega$,
\[
I_0(\Omega)=\bigcup_{\gamma\in\Gamma(\Omega)}L(\gamma)^*,
\]
and let $\acc(\Omega)$ be the set of $x\in C$ for which there is an arc at $x$ in $\Omega$. Note that $I_0(\Omega)$ is
open in $C$ and $\acc(\Omega)\subseteq\overline\Omega$.
\end{definition}

\begin{lemma}[{\cite[Lemma 2]{MMJ}}]\label{lem:mm-lemma2}
\lean{TJK.McMillan.countable_acc_diff_I0, TJK.McMillan.exists_arc_join,
  TJK.McMillan.exists_Ioo_subset_crossArcOpen}
\uses{def:acc, lem:intervals-open, fact:tietze-arc, fact:crosscut}
For every domain $\Omega\subseteq\D\setminus\{0\}$, $\acc(\Omega)\setminus I_0(\Omega)$ is countable.
\end{lemma}
\begin{proof}
We may assume $\acc(\Omega)$ has two or more points. Let $a\in\acc(\Omega)$ and $a'\ne a$ in $\acc(\Omega)$,
with arcs $\gamma,\gamma'$ at $a,a'$ in $\Omega$, and join their other endpoints by a path in $\Omega$. By
Fact~\ref{fact:tietze-arc} the resulting path from $a$ to $a'$ contains a simple arc $\delta$ from $a$ to $a'$;
it is a crosscut with $\delta\cap\D\subseteq\Omega$ (so $0\notin\delta$). Write $a=e^{i\theta}$,
$a'=e^{i(\theta+\ell)}$, $0<\ell<2\pi$. By Theorem~\ref{fact:crosscut}, $L(\delta)^*$ is one of the two open arcs
$e^{i(s,t)}$ with $(s,t)=(\theta,\theta+\ell)$ or $(\theta+\ell,\theta+2\pi)$, and $a\in e^{i[s,t]}$. So
$\acc(\Omega)$ lies in a union of closed arcs $e^{i[s,t]}$ whose open arcs $e^{i(s,t)}$ lie in $I_0(\Omega)$, and
Lemma~\ref{lem:intervals-open} shows that $\acc(\Omega)\setminus I_0(\Omega)$ is countable.
\end{proof}

\begin{lemma}[{\cite[Lemma 3]{MMJ}}]\label{lem:mm-lemma3}
\lean{TJK.McMillan.inter_nonempty, TJK.McMillan.inter_nonempty_of_split, TJK.McMillan.exists_split}
\uses{def:acc, fact:crosscut}
Let $\Omega_1,\Omega_2\subseteq\D\setminus\{0\}$ be connected. If the sets
$I_0(\Omega_1)\cap\overline{\acc(\Omega_1)}$ and $I_0(\Omega_2)\cap\overline{\acc(\Omega_2)}$ have a common
point $a$, then $\Omega_1\cap\Omega_2\ne\emptyset$.
\end{lemma}
\begin{proof}
Suppose $\Omega_1\cap\Omega_2=\emptyset$. Choose $\gamma_i\in\Gamma(\Omega_i)$ with $a\in L(\gamma_i)^*$, and let
$U_i,V_i$ be the components of $\D\setminus\gamma_i$, with $0\in U_i$; so $a\notin\overline{U_i}$, while
$a\in\overline{\acc(\Omega_i)}\subseteq\overline{\Omega_i}$. The sets $\tau_i=\gamma_i\cap\D$ are nonempty, connected,
contained in $\Omega_i$ and in $\overline{U_i}$ (Theorem~\ref{fact:crosscut}), and $\tau_1\cap\gamma_2=\emptyset=
\tau_2\cap\gamma_1$. Hence each $\tau_i$ lies in one component of the other complement.

If $\tau_2\subseteq U_1$: the connected set $\Omega_2$ misses $\gamma_1$ (as $\gamma_1\cap\D\subseteq\Omega_1$), so
it lies in $U_1$ or $V_1$; it meets $U_1$ (it contains $\tau_2$), so $\Omega_2\subseteq U_1$ and
$a\in\overline{\Omega_2}\subseteq\overline{U_1}$, a contradiction. Symmetrically $\tau_1\not\subseteq U_2$.

So $\tau_1\subseteq V_2$ and $\tau_2\subseteq V_1$. As $\tau_1\subseteq\overline{U_1}$ and $V_2$ is open, $U_1$ meets
$V_2$; also $0\in U_1\cap U_2$. The connected set $U_1$ meets both components of $\D\setminus\gamma_2$, so it meets
$\gamma_2\cap\D=\tau_2\subseteq V_1$, which is absurd.
\end{proof}

\begin{theorem}[McMillan {\cite[Theorem 5]{McM}}; {\cite[Theorem 1]{MMJ}}]\label{thm:mcmillan}
\lean{TJK.McMillan.isFσδ_curvConvSet, TJK.McMillan.isFσδ_curvConvSet',
  TJK.McMillan.isFσδ_of_skeleton, TJK.McMillan.exists_path_of_seq,
  TJK.McMillan.exists_arcAt_of_chain}
\uses{lem:mm-lemma2, lem:mm-lemma3, cor:path-arc, lem:subarc, lem:open-components, lem:borel-basic}
Let $K$ be a complete separable metric space and $f:\D\to K$ continuous. Then $\CC(f)$ is an $\Fsd$
subset of $C$.
\end{theorem}
\begin{proof}
Let $\{p_k\}$ be dense in $K$ and $\{c_m\}$ dense in $\C$. Put
$Q(n,m)=S(1/(n+1),c_m)\cap\D\setminus\{0\}$, and let $U(n,m,k,l)$, $l\in\N$, be the component of
$O(n,m,k)=f^{-1}(S(2^{-n},p_k))\cap Q(n,m)$ containing $c_l$ (empty if $c_l\notin O(n,m,k)$); every component of the
open set $O(n,m,k)$ is open and nonempty, so it is one of the $U(n,m,k,l)$ (Lemma~\ref{lem:open-components}). Put
$A(n,m,k,l)=\acc U(n,m,k,l)$ and
\[
Y=\bigcap_n\bigcup_{m,k,l}\Bigl(I_0(U(n,m,k,l))\cap\overline{A(n,m,k,l)}\Bigr).
\]
$I_0(\cdot)$ is open in $C$, so $Y$ is $\Fsd$.

\emph{$Y\subseteq\CC(f)$.} Let $y\in Y$ and choose $U_n=U(n,m[n],k[n],l[n])$ with
$y\in I_0(U_n)\cap\overline{\acc U_n}$. By Lemma~\ref{lem:mm-lemma3}, $U_n\cap U_{n+1}$ contains a
point $z_n$; join $z_n$ to $z_{n+1}$ by a path $\gamma_n\subseteq U_{n+1}$. Since
$y\in\overline{U_{n+1}}\subseteq\overline{S(1/(n+2),c_{m[n+1]})}$, $\gamma_n$ lies within $2/(n+2)$ of $y$. So
the paths $\gamma_n$, traversed on $[n,n+1]$ and reparametrized by $t\mapsto(1-t)^{-1}-1$, followed by $y$ at $t=1$,
form a path $g$ with $g(1)=y$, $g([0,1))\subseteq\D$, and $g(t)\in\bigcup_{n\ge N}\gamma_n$ for $t$ near $1$. As
$S(2^{-n},p_{k[n]})$ and $S(2^{-n-1},p_{k[n+1]})$ meet (at $f(z_n)$), $\rho(p_{k[n]},p_{k[n+1]})<2^{-n+1}$, so
$p_{k[n]}\to p$ for some $p\in K$; and $f(\gamma_n)\subseteq S(2^{-n-1},p_{k[n+1]})$, so $f(g(t))\to p$ as $t\to1$.
By Corollary~\ref{cor:path-arc}, $y\in\CC(f)$.

\emph{$\CC(f)\subseteq\bigcap_n\bigcup_{m,k,l}A(n,m,k,l)$.} Let $f\to p'$ along an arc $\gamma_0$ at $x$.
Given $n$, pick $k$ with $\rho(p',p_k)<2^{-n}$ and $m$ with $\abs{x-c_m}<1/(2n+2)$, so that
$S(r,x)\cap\D\subseteq Q(n,m)$ for $r=\min(1/(2n+2),1/2)$. By Lemma~\ref{lem:subarc} a subarc $\gamma_0'$ at $x$ has
$\gamma_0'\setminus\{x\}\subseteq O(n,m,k)$; being connected, it lies in one component, so $x\in A(n,m,k,l)$
for some $l$.

\emph{$\CC(f)\setminus Y$ is countable.} If $x\in\CC(f)\setminus Y$, there is $n$ with
$x\in A(n,m,k,l)$ for some $m,k,l$ but $x\notin I_0(U(n,m,k,l))\cap\overline{A(n,m,k,l)}$, so
$x\in A(n,m,k,l)\setminus I_0(U(n,m,k,l))$, which is countable by Lemma~\ref{lem:mm-lemma2}.

So $\CC(f)$ is the union of an $\Fsd$ set and a countable set, hence $\Fsd$. (In Lean the three steps are
packaged once, for any countable family of open sets $O(n,i)\subseteq\D\setminus\{0\}$, and reused for
Theorem~\ref{thm:mm-thm2}.)
\end{proof}

\begin{theorem}[{\cite[Theorem 2]{MMJ}}]\label{thm:mm-thm2}
\lean{TJK.McMillan.isFσδ_clusterSet_subset, TJK.McMillan.isFσδ_clusterSet_subset',
  TJK.McMillan.mem_map_of_forall_mapClusterPt}
\uses{lem:mm-lemma2, lem:mm-lemma3, cor:path-arc, lem:subarc, lem:cluster, lem:open-components,
  lem:borel-basic, fact:hyperspace}
Let $K$ be a compact metric space, $\mathcal C(K)$ the space of nonempty closed subsets of $K$ with the
Hausdorff metric $\bar\rho$, $\mathcal E\subseteq\mathcal C(K)$ closed, and $f:\D\to K$ continuous. Then
\[
X_{\mathcal E}=\{x\in C:\text{there are an arc $\gamma$ at $x$ and $E\in\mathcal E$ with }C(f,\gamma)\subseteq E\}
\]
is an $\Fsd$ set.
\end{theorem}
\begin{proof}
If $\mathcal E=\emptyset$ there is nothing to prove. $\mathcal E$ is compact (Fact~\ref{fact:hyperspace}), hence has a dense sequence
$P(k)$. For $\eps>0$ let $\mathcal S(\eps,E)=\{a:\dist(a,E)<\eps\}$, which is open, and note that
$\dist(a,F)\le\dist(a,E)+\bar\rho(E,F)$. Define $Q(n,m)$ as before, $U(n,m,k,l)$ as the
components of $f^{-1}(\mathcal S(1/(n+1),P(k)))\cap Q(n,m)$, $A(n,m,k,l)=\acc U(n,m,k,l)$, and $Y$ by the same
formula; $Y$ is $\Fsd$.

\emph{$Y\subseteq X_{\mathcal E}$.} With $U_n$ as before, compactness gives $P\in\mathcal E$ and
$n_1<n_2<\cdots$ with $P(k[n_j])\to P$. The sets $U_{n_j}$ satisfy the hypotheses used above, so as in
Theorem~\ref{thm:mcmillan} we get a path $g$ ending at $y$, with $g(t)\in\bigcup_{j\ge J}U_{n_j}$ for $t$ near $1$,
and (Corollary~\ref{cor:path-arc}) an arc $\gamma$ at $y$ whose cluster set is contained in that of $f\circ g$ at $1$.
Let $p$ be such a cluster point. Given $\eps>0$, for $j$ large $2/(n_j+1)<\eps/2$ and
$\bar\rho(P(k[n_j]),P)<\eps/2$, so $\dist(f(g(t)),P)\le\eps$ for $t$ near $1$, and hence $\dist(p,P)\le\eps$. As $P$
is closed, $p\in P$. So $C(f,\gamma)\subseteq P$.

\emph{$X_{\mathcal E}\subseteq\bigcap_n\bigcup_{m,k,l}A(n,m,k,l)$.} Let $C(f,\gamma_0)\subseteq P_0\in\mathcal E$.
Given $n$, choose $k$ with $\bar\rho(P_0,P(k))<1/(n+1)$, so $C(f,\gamma_0)\subseteq\mathcal S(1/(n+1),P(k))=:O$. Some
subarc of $\gamma_0$ at $x$ is mapped into $O$: otherwise the filter of $f$ along $\gamma_0$ meets the compact
set $K\setminus O$ and has a cluster point there, which would lie in $C(f,\gamma_0)\setminus O$. Conclude as in
Theorem~\ref{thm:mcmillan}; the countability of $X_{\mathcal E}\setminus Y$ is proved in the same way.
\end{proof}

\begin{corollary}[{\cite[Corollaries 1--4]{MMJ}}]\label{cor:mm-corollaries}
\lean{TJK.McMillan.cor1, TJK.McMillan.cor2, TJK.McMillan.cor3, TJK.McMillan.cor4,
  TJK.McMillan.isClosed_subset_greatCircle}
\uses{thm:mm-thm2, lem:cluster}
Let $f:\D\to\Sph$ be continuous. Each of the following sets is $\Fsd$ ((1)--(3) hold for any compact metric
target):
\begin{enumerate}[label=(\arabic*)]
\item (McMillan) $\{x:\exists\gamma\text{ at }x,\ \exists p\in E,\ f\to p\text{ along }\gamma\}$, for closed
  $E\subseteq\Sph$;
\item $\{x:\exists\gamma\text{ at }x,\ \diam C(f,\gamma)\le d\}$, for $d\ge0$;
\item $\{x:\exists\gamma\text{ at }x,\ C(f,\gamma)\subseteq E\}$, for closed $E\subseteq\Sph$;
\item $\{x:\exists\gamma\text{ at }x,\ C(f,\gamma)\text{ is an arc of a great circle}\}$, where (since
  $C(f,\gamma)$ is closed and connected) an ``arc'' may be a point, a closed arc, or a whole great circle; that is,
  $C(f,\gamma)$ is contained in a great circle.
\end{enumerate}
\end{corollary}
\begin{proof}
Apply Theorem~\ref{thm:mm-thm2} with $\mathcal E$ equal to: (1) $\{\{p\}:p\in E\}$, an isometric copy of
the compact set $E$, using that $C(f,\gamma)\subseteq\{p\}$ iff $f\to p$ along $\gamma$ (Lemma~\ref{lem:cluster});
(2) $\{F:\diam F\le d\}$, closed because if $\bar\rho(F,F')<\eps$ then every two points of $F$ are within
$2\eps$ of two points of $F'$, and using $E=C(f,\gamma)$ for the converse; (3) $\{E\}$ (or note the set is empty if
$E=\emptyset$); (4) the sets contained in some great circle $\{v:\langle v,u\rangle=0\}$: if $F_n\to F$ with
$F_n\subseteq\{\langle\cdot,u_n\rangle=0\}$, pass to a subsequence $u_n\to u$; each $v\in F$ is a limit of
$v_n\in F_n$, so $\langle v,u\rangle=\lim\langle v_n,u_n\rangle=0$. With $E=C(f,\gamma)$ for the converse; the
connectedness of $C(f,\gamma)$ (Lemma~\ref{lem:cluster}) justifies the reading of ``arc''.
\end{proof}

\begin{theorem}[McMillan in the half-plane; {\cite[Theorem 3]{Diss}}]\label{thm:mcmillan-H}
\lean{TJK.McMillan.isFσδ_curvConvSet_upperHalfPlane}
\uses{thm:mcmillan, lem:transfer, lem:charts, lem:borel-basic}
If $K$ is a complete separable metric space and $f:\Hp\to K$ is continuous, then $\CC(f)\subseteq\R$ is
$\Fsd$.
\end{theorem}
\begin{proof}
By Lemma~\ref{lem:transfer} with $\Psi=\kappa$ (Lemma~\ref{lem:charts}(a)),
$\CC(f)=\R\cap\kappa^{-1}(\CC(f\circ\kappa^{-1}))$, and the right side is $\Fsd$ by
Theorem~\ref{thm:mcmillan} and Lemma~\ref{lem:borel-basic} (preimage under the continuous map $\kappa|_\R$, then
an $\Fsd$ subset of the closed set $\R$).
\end{proof}

\subsection*{The half-plane lemmas of the dissertation}
For a domain $\Omega\subseteq\Hp$ and a crosscut $\gamma$ of $\Hp$ with endpoints $c<d$ let
$L(\gamma)=[c,d]$, the trace on $\R$ of the closure of the bounded component of $\Hp\setminus\gamma$; let
$I(\Omega)=\bigcup L(\gamma)^*$ over crosscuts with $\gamma\cap\Hp\subseteq\Omega$, and $\acc(\Omega)$ the
points of $\R$ accessible by arcs in $\Omega$. These lemmas are the half-plane forms of
Lemmas~\ref{lem:mm-lemma2} and~\ref{lem:mm-lemma3}; they are not used elsewhere.

\begin{lemma}[{\cite[Lemma 7]{Diss}}]\label{lem:th-lemma7}
\lean{TJK.McMillan.halfPlane_I_eq, TJK.McMillan.IH, TJK.McMillan.accH}
\uses{lem:crosscut-H, fact:tietze-arc}
If $\acc(\Omega)\ne\emptyset$, $a=\inf\acc(\Omega)$ and $b=\sup\acc(\Omega)$, then $I(\Omega)=(a,b)$; that is,
$x\in I(\Omega)$ iff there are $c,d\in\acc(\Omega)$ with $c<x<d$.
\end{lemma}
\begin{proof}
If $x\in L(\gamma)^*=(c,d)$ then $c,d\in\acc(\Omega)$ (initial halves of $\gamma$ from either end are arcs at $c$,
$d$ in $\Omega$), so $a\le c<x<d\le b$. Conversely, for $a<x<b$ pick $c',d'\in\acc(\Omega)$ with $c'<x<d'$, and join
them by a crosscut through $\Omega$ as in Lemma~\ref{lem:mm-lemma2}.
\end{proof}

\begin{lemma}[{\cite[Lemma 8]{Diss}}]\label{lem:th-lemma8}
\lean{TJK.McMillan.halfPlane_inter_nonempty}
\uses{lem:crosscut-H, lem:mm-lemma3}
If $\Omega_1,\Omega_2\subseteq\Hp$ are connected and $I(\Omega_1)\cap\overline{\acc(\Omega_1)}$ and
$I(\Omega_2)\cap\overline{\acc(\Omega_2)}$ meet, then $\Omega_1\cap\Omega_2\ne\emptyset$.
\end{lemma}
\begin{proof}
As in Lemma~\ref{lem:mm-lemma3}, with $U_i$ the unbounded component of $\Hp\setminus\gamma_i$ (any two
unbounded components meet, since each contains a half-plane $\{\operatorname{Im}z>t\}$) and
Lemma~\ref{lem:crosscut-H} in place of Theorem~\ref{fact:crosscut}.
\end{proof}
